\chapter{Literature Review}
\label{chap:litreview}

\section{Introduction}

Wildfires are increasingly recognized as complex socioecological hazards whose consequences extend far beyond fire occurrence itself. Recent literature shows that machine learning, deep learning, remote sensing, geospatial data fusion, and multimodal artificial intelligence have substantially improved the ability to predict wildfire ignition, susceptibility, spread, fire danger, active-fire detection, and short-term fire behavior. At the same time, a smaller but rapidly growing body of work connects wildfire prediction to broader consequence assessment, including suppression expenditures, direct asset damage, expected annual loss, supply-chain disruption, smoke-related health burdens, and ecological degradation. This chapter reviews influential and recent literature across these connected areas and positions the present thesis at the intersection of wildfire-risk prediction and wildfire-loss assessment.

The review is organized around several connected themes. First, it clarifies the conceptual foundations of wildfire risk, hazard, exposure, vulnerability, resilience, and loss. Second, it discusses foundational fire-behavior, fire-regime, and wildfire-climate studies that motivate the use of environmental predictors. Third, it reviews data foundations, including remote sensing, meteorological datasets, land-cover products, and wildfire benchmark datasets. Fourth, it examines machine-learning and deep-learning approaches for wildfire prediction and monitoring. Fifth, it reviews the literature that moves from wildfire prediction to economic loss assessment, including expected annual loss, suppression costs, property damage, supply-chain disruption, economy-wide impacts, and smoke-related mortality. Sixth, it discusses agricultural and natural-hazard loss modeling as a methodological bridge. Seventh, it reviews spatial transferability and geographic validation, ecological losses, and emerging LLM/foundation-model approaches. Finally, it identifies the research gap addressed by this thesis.

\section{Wildfire Risk, Hazard, Vulnerability, and Loss}

A recurring issue in the wildfire literature is the inconsistent use of the terms \emph{risk}, \emph{hazard}, \emph{danger}, \emph{exposure}, \emph{vulnerability}, \emph{resilience}, and \emph{loss}. Recent wildfire-risk reviews emphasize that wildfire risk cannot be reduced to fire occurrence alone, because the impact of a wildfire depends not only on where and when fire occurs, but also on what assets, populations, infrastructure, and ecosystems are exposed and how susceptible they are to damage \citep{xu2024survey,xu2025deeplearning,jodhani2026ai,caron2025ai,ejaz2025mlpipeline}. In this sense, wildfire hazard generally refers to the likelihood, intensity, or physical potential of fire occurrence, while vulnerability captures the susceptibility of exposed systems to damage and their capacity to recover \citep{arrogantefunes2024ecological,asif2026supranational}.

This distinction is central to the present thesis because much of the machine-learning wildfire literature remains hazard-centric. Many studies stop at susceptibility maps, ignition probability, spread forecasts, or fire-danger scores, while fewer translate predictive outputs into explicit estimates of economic or ecological loss \citep{asif2026supranational,caron2025ai}. Wildfire-loss assessment therefore requires a broader framework in which hazard, exposure, vulnerability, and resilience are explicitly combined to estimate expected consequences \citep{asif2026supranational,ma2022supplychain,monge2023probabilistic,thompson2011risk,calkin2014risk}.

Influential fire-regime literature also shows that wildfire risk is shaped by the interaction of climate, vegetation, land use, and human systems. Bowman et al.\ describe fire as a central Earth-system process, while later work emphasizes the human dimension of global fire regimes, showing that fire is not only a biophysical process but also a product of land management, settlement patterns, ignition sources, and suppression practices \citep{bowman2009fire,bowman2011human}. Moritz et al.\ argue that societies need to learn to coexist with wildfire by shifting from a purely suppression-oriented paradigm to a risk-reduction and resilience perspective \citep{moritz2014learning}. These studies support the view that wildfire loss cannot be explained by fire occurrence alone; it depends on the coupled human--environment system in which fire occurs.

A similar conceptual separation is needed when distinguishing economic and ecological losses. Economic losses include suppression expenditures, damage to property and infrastructure, indirect business interruptions, supply-chain disruptions, public-health costs linked to smoke exposure, and broader macroeconomic impacts \citep{huang2024suppression,wang2021economicfootprint,monge2023probabilistic,law2025anthropogenic,qiu2025wildfiresmoke}. Ecological losses include biodiversity reduction, vegetation degradation, loss of ecosystem functions, reduced ecosystem-service capacity, reduced resilience, and long recovery periods after severe fires \citep{arrogantefunes2024ecological,rocesdiaz2022ecosystemservices,pacheco2023ecosystemservices,taboada2021ecosystemdelivery}. A complete wildfire-loss framework should therefore not rely on burned area alone as a proxy for impact, but should account for the value of exposed systems, the severity of damage, and recovery dynamics.

\section{Foundational Fire-Behavior, Climate, and Fire-Regime Literature}

Wildfire modeling has long been supported by physical and empirical fire-behavior models. Rothermel's surface fire-spread model remains one of the foundational contributions for predicting fire spread in wildland fuels \citep{rothermel1972spread}. The Canadian Forest Fire Weather Index system provides another influential operational framework for translating weather variables into fire-danger indicators \citep{vanwagner1987fwi}. FARSITE extended fire-behavior modeling into spatial simulation, while fuel-model references by Anderson and Scott and Burgan formalized important fuel categories used in fire-behavior estimation \citep{finney1998farsite,anderson1982fuelmodels,scott2005fuelmodels}. BehavePlus and related fire-behavior systems further demonstrate the long-standing operational importance of structured fuel, weather, and topographic inputs \citep{andrews2014behaveplus}.

Although this thesis uses machine learning rather than a physics-based fire-spread simulator, these foundational works are still important because they clarify why terrain, fuels, wind, moisture, and weather are central wildfire predictors. They also provide conceptual continuity between classical fire-science variables and modern data-driven feature engineering. Earlier probability-based and GIS/remote-sensing wildfire-risk frameworks similarly demonstrate that fire risk has long been modeled as a combination of hazard, environmental conditions, exposure, and values at risk \citep{preisler2004forecasting,chuvieco2010framework,ager2010wildfire,thompson2011risk}.

A large body of influential wildfire-climate research links changing fire regimes to warming temperatures, increasing atmospheric dryness, and longer fire-weather seasons. Westerling et al.\ showed that warming and earlier spring conditions contributed to increased western U.S. forest wildfire activity, establishing one of the key empirical links between climate change and wildfire activity \citep{westerling2006warming}. Jolly et al.\ demonstrated that global fire-weather seasons lengthened across many regions between 1979 and 2013, indicating that climate change can expand the temporal window in which large fires are possible \citep{jolly2015climate}. Flannigan et al.\ and Moritz et al.\ further discuss how changing climate can alter global fire activity, while Barbero et al.\ show increasing potential for very large fires in the contiguous United States \citep{flannigan2009climate,moritz2012climate,barbero2015climate}.

Environmental controls on wildfire operate across multiple spatial and temporal scales. Parisien and Moritz analyze environmental controls on wildfire distribution at multiple scales, while Littell et al.\ connect climate variability to area burned across western U.S. ecoprovinces \citep{parisien2009environmental,littell2009climate}. Krawchuk et al.\ introduce the broader concept of global pyrogeography, and Abatzoglou et al.\ show that climate--fire relationships vary across regions \citep{krawchuk2009global,abatzoglou2018global}. Human influence is also central: Syphard et al.\ examine human influence on California fire regimes, and Balch et al.\ show that human-started wildfires expand the fire niche across the United States \citep{syphard2007human,balch2017human}. In tropical forest systems, Arag{\~a}o and Shimabukuro show that fire incidence in Amazonian forests is closely connected to land-use change and has important implications for forest conservation and REDD policy \citep{aragao2010incidence}. Dennison et al.\ document large-wildfire trends in the western United States, reinforcing the importance of long-term fire-regime changes \citep{dennison2014large}.

Recent work highlights atmospheric aridity and moisture availability. Abatzoglou and Williams show that anthropogenic climate change increased fuel aridity across western U.S. forests and contributed substantially to increased forest fire area \citep{abatzoglou2016anthropogenic}. Williams et al.\ further show that increased atmospheric aridity played a major role in observed increases in California wildfire activity \citep{williams2019observed}. Trugman et al.\ emphasize the role of moisture availability and timing, while Holden et al.\ show that soil moisture and plant hydraulic stress can be strong predictors of fire-spread potential \citep{trugman2018flash,holden2025soilmoisture}. These studies provide important physical justification for including climate variables such as precipitation, maximum temperature, wind speed, and vapor pressure deficit in wildfire-risk and wildfire-loss models.

Global fire and emissions studies also provide context for wildfire impacts. Global fire-emissions datasets and fire radiative power methods show how remote sensing can quantify combustion, emissions, and biomass burning activity at large scales \citep{randerson2012gfed,van_der_werf2017global,wooster2005fire}. The State of Wildfires 2024--2025 report provides a recent global assessment of extreme wildfire episodes, exposure, emissions, and climate drivers \citep{kelley2025state}. Together, foundational fire-behavior, climate, and fire-regime studies provide the scientific basis for the environmental component of the present thesis.

\section{Data Foundations for Wildfire Modeling}

The recent wildfire literature is strongly data-driven. Review papers consistently show that wildfire-risk models are built from a combination of remote-sensing imagery, meteorological and climate data, vegetation and fuels information, topography, hydrology, historical fire records, and socioeconomic variables \citep{jain2020review,xu2025deeplearning,xu2024survey,ejaz2025mlpipeline,jodhani2026ai,liu2025advancements}. Satellite platforms such as Landsat, Sentinel, MODIS, VIIRS, GOES, LiDAR, SAR, and UAV-based systems have become central to wildfire research because they enable broad spatial coverage, repeated observations, and the derivation of key indices such as NDVI, NDWI, NBR, land surface temperature, canopy structure, burn scars, and fuel moisture \citep{jodhani2026ai,xu2025deeplearning}.

Remote-sensing fire products are especially important for active fire detection, burned-area mapping, and burn-severity assessment. The Monitoring Trends in Burn Severity (MTBS) project provides an influential reference for tracking burn severity over time \citep{eidenshink2007mtbs}. The MODIS active fire product and the VIIRS 375 m active fire product are widely used for active fire detection and near-real-time fire monitoring \citep{giglio2016modis,schroeder2014viirs}. Burn severity and spectral indices are also important for post-fire impact analysis: Key and Benson formalized the Composite Burn Index and normalized burn ratio approach, while Miller and Thode developed the relative differenced normalized burn ratio for heterogeneous landscapes \citep{key2006firemon,miller2007severity}.

Recent benchmark-oriented work highlights that data-source selection is itself a major methodological decision. Xu et al.\ introduce Sen2Fire, a Sentinel-based wildfire detection benchmark that combines Sentinel-2 multispectral imagery with Sentinel-5P aerosol information and evaluates the role of spectral bands and indices such as NBR and NDVI \citep{xu2024sen2fire}. Karlsson et al.\ compare MODIS and VIIRS products for next-day wildfire predictability and show that the choice of satellite fire-mask product can substantially affect learning behavior and prediction quality \citep{karlsson2025modisviirs}. Huot et al.\ provide a machine-learning dataset for next-day wildfire spread prediction from remote-sensing data, illustrating how wildfire modeling is moving toward benchmarked spatiotemporal learning tasks \citep{huot2022nextday}.

Environmental datasets also play a major role. The gridMET dataset provides gridded daily surface meteorological variables designed for ecological and environmental applications \citep{abatzoglou2013gridmet}. The National Land Cover Database provides Landsat-based land-cover products at 30 m resolution, making it suitable for deriving county-level land-cover fractions such as forest, shrub, grass, and developed land \citep{dewitz2021nlcd}. Vegetation and fuel-mapping work further shows why fuel and land-cover information are essential for fire-risk assessment and fire-management planning \citep{keane2001mapping}. These datasets are directly relevant to the present thesis because they support scalable construction of county-level environmental predictors for wildfire expected annual loss modeling.

Despite these advances, the literature identifies several data limitations. Cross-ecosystem transferability remains weak because models trained in one biome can perform poorly in another. Real-time integration remains limited in many studies, which often rely on static or historical datasets rather than dynamically updated data streams. Uncertainty quantification is also underdeveloped, especially for operational systems where probabilistic outputs would often be more useful than single deterministic scores \citep{jodhani2026ai,caron2025ai}. Wildfire datasets also frequently exhibit class imbalance, spatial autocorrelation, and inconsistent labeling across data products \citep{xu2024survey,caron2026multitarget,karlsson2025modisviirs}.

A second challenge is that data sufficient for wildfire prediction are not always sufficient for wildfire-loss estimation. Loss-focused studies require additional exposure and vulnerability layers, such as buildings, transport links, forests, energy infrastructure, demographic profiles, agricultural values, and other economic-value indicators \citep{asif2026supranational,ma2022supplychain,wang2021economicfootprint}. The transition from wildfire prediction to wildfire-loss modeling is therefore both a methodological shift and a data-expansion problem.

\section{Machine Learning and Deep Learning for Wildfire Prediction}

\subsection{Traditional Machine Learning}

Traditional machine-learning models, including logistic regression, random forests, support vector machines, gradient boosting, and related ensemble methods, are widely used for wildfire ignition, susceptibility, and danger mapping \citep{jain2020review,sayad2019predictive,liu2025advancements,xu2024survey,jodhani2026ai,ejaz2025mlpipeline}. These models are attractive because they can capture nonlinear relationships among weather, topography, vegetation, and human variables while remaining more interpretable and computationally accessible than many deep-learning approaches.

Recent reviews suggest that traditional ML methods remain valuable as robust baselines, especially where interpretability, lower computational cost, and limited training data are important \citep{jain2020review,jodhani2026ai,liu2025advancements,ejaz2025mlpipeline}. However, their performance can deteriorate when spatial heterogeneity is high, when labels are imbalanced, or when the task shifts from susceptibility mapping to real-time fire monitoring or complex spatiotemporal forecasting \citep{caron2025ai}.

\subsection{Deep Learning and Spatiotemporal Prediction}

Deep learning has become increasingly important in wildfire research because of its ability to model high-dimensional spatiotemporal data and learn interactions across imagery, fuels, weather, and terrain. Recent review papers group these approaches into image classification, semantic segmentation, time-series prediction, fire spread forecasting, and broader spatiotemporal risk modeling \citep{xu2024survey,xu2025deeplearning,ejaz2025mlpipeline}. Convolutional neural networks, recurrent neural networks, LSTMs, ConvLSTMs, graph neural networks, transformers, and multimodal architectures have been used for ignition likelihood, fire spread, danger estimation, detection, and burned-area progression \citep{liu2025advancements,xu2025deeplearning}. Earlier deep-learning examples such as FireCast demonstrate the potential of deep models for wildfire-spread prediction, while more recent benchmark datasets support more systematic evaluation \citep{radke2019firecast,huot2022nextday}.

The review literature is broadly consistent in its assessment: deep learning improves predictive flexibility and can achieve strong performance, but operational value depends on interpretability, calibration, transferability, data quality, and computational feasibility \citep{jodhani2026ai,liu2025advancements,caron2025ai}. Data-assimilation studies also indicate a shift toward dynamic model--data fusion. For example, Fire-EnSF proposes an ensemble score filter for wildfire spread data assimilation, integrating numerical fire models with observations to improve real-time spread prediction \citep{shi2025fireensf}. Such work shows that wildfire forecasting is moving beyond static mapping toward dynamic and continuously updated systems.

\subsection{Operational Gaps in Wildfire AI}

Although wildfire AI research is growing rapidly, operational deployment still lags behind research output. Caron et al.\ argue that dataset imbalance, weak benchmarking, unsuitable metrics, computational demands, and limited integration into decision-support workflows remain major barriers \citep{caron2025ai}. Their multi-target wildfire-risk proof of concept further emphasizes that operational wildfire decision-making cannot be reduced to a single risk score: it requires simultaneous consideration of meteorological danger, ignition activity, intervention complexity, and resource mobilization \citep{caron2026multitarget}.

These observations are relevant for the present thesis because they indicate that accurate wildfire prediction is necessary but not sufficient. To support decision-making, models must produce interpretable and scalable outputs that can be linked to consequence analysis, including expected annual losses and broader economic or ecological impacts.

\section{From Wildfire Prediction to Economic Loss Assessment}

\subsection{Suppression Costs and Operational Expenditure}

One strand of recent research examines the relationship between wildfire detection, response timing, and suppression costs. Huang and Wichmann use machine learning and double/debiased machine learning to estimate how reporting delays affect suppression expenditures in Alberta, Canada \citep{huang2024suppression}. Their findings show that faster detection can reduce suppression costs, but that suppression savings alone may not fully justify investments in detection. This is important because the value of early detection also includes avoided property losses, ecological impacts, public-health burdens, and indirect economic disruption.

Ardid et al.\ similarly connect forecast skill with economic value in wildfire potential forecasting, demonstrating that prediction quality should be assessed not only statistically but also in terms of decision value \citep{ardid2026forecastskill}. This perspective is important for a loss-oriented thesis because it shifts evaluation from accuracy alone toward the consequences of improved prediction.

\subsection{Expected Annual Loss and Public Risk Datasets}

A more comprehensive approach explicitly links hazard modeling with exposure and vulnerability to estimate losses. The supranational study by Asif et al.\ develops a machine-learning framework for estimating wildfire losses under climate change in Southeastern Europe \citep{asif2026supranational}. Their method combines random-forest-based susceptibility mapping, hazard classification, exposure layers, vulnerability tables, and economic valuation to estimate average annual loss. This study is especially close to the present thesis because it bridges hazard mapping and loss estimation rather than treating them as separate research streams.

For scalable wildfire-loss research, public risk datasets are also important because they provide standardized targets and exposure-related indicators across large regions. The FEMA National Risk Index provides county- and census-tract-level indicators for multiple natural hazards, including wildfire \citep{fema2025nri_technical,fema2025nri_methodology}. The NRI framework combines expected annual loss, social vulnerability, and community resilience to characterize natural-hazard risk. For wildfire, this makes it possible to move from a purely occurrence-based target toward a monetary risk target that reflects expected consequences.

Expected annual loss is useful because it provides a comparable risk measure across counties and hazards. Instead of asking only whether a wildfire may occur, EAL asks what average annual loss may be expected given hazard frequency, exposure, and consequence assumptions. At the same time, public EAL datasets should be interpreted carefully. They are spatially aggregated, depend on embedded assumptions, and may not capture event-level damage mechanisms directly. They are therefore most appropriate as scalable first baselines that should eventually be complemented by more detailed event-level or asset-level data.

\subsection{Property Damage, Economy-Wide Effects, and Supply Chains}

Beyond suppression and EAL frameworks, wildfire-loss modeling increasingly includes direct structural damage, supply-chain disruption, and economy-wide impacts. Jia and Opabola develop an interpretable machine-learning framework for wildfire damage drivers in California, using explainable AI to identify how hazard, exposure, and vulnerability jointly shape damage outcomes \citep{jia2025interpretable}. Raveendran et al.\ propose a flexible semiparametric approach to wildfire-loss modeling, linking actuarial loss modeling with more flexible statistical methods \citep{raveendran2025wildfireloss}.

Wang et al.\ estimate the economic footprint of the 2018 California wildfires and show that total damages extend far beyond direct fire impacts, including health and indirect economic effects \citep{wang2021economicfootprint}. Monge et al.\ quantify economy-wide wildfire impacts under changing wildfire climate conditions in a regional rural landscape \citep{monge2023probabilistic}. Ma et al.\ extend the loss perspective to supply-chain networks, where wildfire affects physical assets, system performance, transportation times, unmet demand, and total supply-chain costs \citep{ma2022supplychain}. Together, these studies show that wildfire-loss research is moving beyond burned area and direct property damage toward cascading, functional, and economy-wide consequence analysis.

\subsection{Smoke-Related Health and Mortality Burdens}

A rapidly developing component of wildfire-loss literature concerns smoke-related mortality and health damages. Earlier health-impact reviews and global mortality studies show that wildfire smoke exposure can cause substantial health burdens across regions \citep{johnston2012global,reid2016critical}. Burke et al.\ discuss the changing risk and burden of wildfire in the United States, connecting smoke and exposure to broader societal consequences \citep{burke2021smoke}. More recent studies quantify the contribution of anthropogenic climate change to wildfire particulate matter and related mortality, and estimate future wildfire smoke exposure and mortality burdens under climate change \citep{law2025anthropogenic,qiu2025wildfiresmoke}.

This strengthens the argument that wildfire economic losses are not limited to fire-front impacts. They propagate through public health, regional healthcare burdens, long-range smoke exposure, and productivity effects. For a thesis focused on wildfire losses, smoke-health literature broadens the definition of consequence and highlights the need for integrated models that consider both direct and indirect impacts.

\section{Methodological Lessons from Agricultural and Natural-Hazard Loss Modeling}

Agricultural-loss studies provide useful analogies for wildfire-loss research because they explicitly connect environmental hazards with measurable impacts on production, asset value, and recovery. In contrast to many wildfire studies that stop at susceptibility or hazard mapping, this literature often estimates the consequences of hazard exposure. Lateef et al.\ combine optical and SAR imagery with participatory mapping to model flooded crop damage, loss, and recovery, thereby linking physical hazard detection with post-disaster consequence assessment \citep{lateef2025floodedcrop}. Dang et al.\ construct a crop flood damage assessment index for rapid post-disaster impact estimation \citep{dang2024cropflood}.

This direction is reinforced by studies connecting climate extremes to broader agricultural and socioeconomic losses. Zafar et al.\ estimate cropland and infrastructure damage caused by flooding through remote sensing \citep{zafar2024flooddamages}, while Shi et al.\ quantify crop yield and production responses to drought and flood disasters across China \citep{shi2021cropyield}. M\"uller et al.\ show that Sentinel-2 data can quantify drought-related crop-yield impacts at local and regional scales \citep{mueller2025droughtsweden}, and Zhu et al.\ integrate crop, groundwater, and economic dynamics to analyze adaptation and economic outcomes under climate stress \citep{zhu2025ogallala}. Although these studies are centered on agriculture rather than wildfire, they offer transferable methodological lessons: hazard signals must be connected to exposed systems, quantified consequences, and recovery dynamics if risk modeling is to become genuinely loss-oriented.

\section{Environmental Drivers of Wildfire Risk and Loss}

The environmental component of wildfire risk has been extensively documented in climate and fire literature. Variables such as precipitation, vapor pressure deficit, temperature, wind, vegetation structure, fuel type, soil moisture, elevation, and slope influence fuel moisture, fire behavior, ignition potential, spread, and suppression difficulty. These variables are therefore useful not only for predicting fire occurrence but also for loss-oriented wildfire-risk modeling, because they influence where damaging fires are more likely to occur and how severe those events may become.

VPD is particularly important because it captures atmospheric drying demand and is closely related to fuel aridity. Abatzoglou and Williams show that anthropogenic climate change increased fuel aridity and contributed to increased forest fire area across the western United States \citep{abatzoglou2016anthropogenic}. Williams et al.\ further show that atmospheric aridity played a major role in observed increases in California wildfire activity \citep{williams2019observed}. Holden et al.\ add that soil moisture and plant hydraulic stress can be strong predictors of forest fire spread potential, highlighting the importance of vegetation water stress and hydroclimate variables \citep{holden2025soilmoisture}. Beltrán-Marcos et al.\ show that extreme burned area and severity can be driven by atmospheric conditions such as high VPD, strong winds, drought, and heatwave periods \citep{beltranmarcos2026drivers}.

Precipitation is also central because it affects both fuel moisture and vegetation growth. In some settings, low precipitation directly increases dryness, while in others antecedent precipitation increases biomass and later fuel loading. Land-cover variables such as forest, shrub, grass, and developed fractions characterize the exposed landscape and potential fuel structure. Topographic variables such as elevation and slope influence climate gradients and fire behavior, while wind affects spread rate and suppression difficulty. For this reason, geospatial datasets such as gridMET and NLCD are useful for constructing scalable environmental predictors \citep{abatzoglou2013gridmet,dewitz2021nlcd}.

For the present thesis, this literature is directly relevant because the experimental results show that environmental variables substantially improve county-level wildfire EAL prediction. The literature provides scientific explanation for why predictors such as warm-season precipitation, VPD, shrub fraction, elevation, slope, and wind speed are important in the combined model.

\section{Spatial Transferability and Geographic Validation}

A major challenge in wildfire modeling is spatial transferability. Models trained and tested using random splits may appear accurate because nearby or environmentally similar locations can appear in both training and test sets. However, this can overestimate performance when the intended use case involves prediction in new regions or under different environmental conditions. This issue is widely discussed in spatial ecology, environmental modeling, and geospatial machine learning.

Roberts et al.\ argue that blocked or structured cross-validation is often more appropriate than random cross-validation when data have spatial, temporal, hierarchical, or phylogenetic structure \citep{roberts2017crossvalidation}. Valavi et al.\ introduce the blockCV framework for generating spatially or environmentally separated folds in species distribution modeling, emphasizing that random validation can underestimate prediction error in spatially structured data \citep{valavi2019blockcv}. Meyer and Pebesma discuss the problem of predicting into unknown space and propose the concept of area of applicability, where model predictions are more reliable when new data fall within the environmental conditions represented in the training set \citep{meyer2021aoa}. Wang et al.\ further develop spatial+ cross-validation ideas for geospatial problems, reinforcing the need to evaluate models under spatially realistic validation settings \citep{wang2023spatialplus}.

These ideas are important for wildfire-risk modeling because fire regimes, vegetation types, climate gradients, suppression practices, and exposure patterns vary strongly across regions. A model that performs well under random validation may be less reliable when evaluated across held-out states or regions. For this reason, the present thesis uses state-wise grouped validation as a stricter test of geographic generalization. This validation design is not a replacement for all spatial validation methods, but it provides a more demanding benchmark than random splitting and helps assess whether the model captures transferable wildfire-risk signals.

\section{Ecological Losses, Vulnerability, and Recovery}

Compared with economic-loss modeling, ecological-loss assessment remains difficult because ecosystem damage cannot be fully represented by a single monetary measure. Ecological losses may involve biodiversity decline, vegetation degradation, soil degradation, hydrological changes, reduced habitat quality, loss of ecosystem-service capacity, and delayed recovery after severe fire. The European ecological-vulnerability framework proposed by Arrogante-Funes et al.\ is particularly relevant because it models vulnerability through ecological value before fire, coping capacity during fire, and post-fire recovery time \citep{arrogantefunes2024ecological}. Instead of treating burned area as the full proxy for ecological impact, this framework emphasizes that ecological consequences depend on ecosystem value, resistance, coping capacity, and recovery.

\subsection{Ecological Vulnerability and Ecosystem Services}

Recent ecosystem-service studies provide a broader way to interpret ecological wildfire losses. Roces-D\'iaz et al.\ synthesize global evidence on fire effects on ecosystem services in forests and woodlands, showing that fire effects vary across service types, ecosystem contexts, and fire characteristics \citep{rocesdiaz2022ecosystemservices}. Pacheco et al.\ characterize wildfire impacts on ecosystem services in Portugal by triangulating scientific literature, governmental reports, and expert perceptions, demonstrating that wildfire impacts are not limited to vegetation loss but also affect water regulation, erosion control, recreation, cultural values, and provisioning services \citep{pacheco2023ecosystemservices}. Taboada et al.\ show that wildfire can reduce ecosystem-service delivery in maritime pine-dominated forests, including provisioning and cultural services \citep{taboada2021ecosystemdelivery}. Misseyanni et al.\ further review forest-fire impacts on ecosystem services in Greece, emphasizing negative effects on hydrological regulation, erosion control, cultural services, climate regulation, and provisioning services \citep{misseyanni2025ecosystemservices}.

These studies are important for the present thesis because they show that ecological losses are multidimensional. A wildfire may produce direct vegetation damage, but its consequences also extend to ecosystem functions and services that support human and natural systems. Therefore, ecological loss should be treated as a separate but connected component of wildfire-loss assessment rather than as a secondary effect of economic damage.

\subsection{Post-Fire Vegetation Recovery and Remote Sensing}

Post-fire recovery is a central dimension of ecological loss because two areas with similar burned area may differ greatly in recovery speed and long-term ecological degradation. P\'erez-Cabello et al.\ review remote-sensing techniques for assessing post-fire vegetation recovery and show that satellite time series, spectral indices, and multi-temporal analysis can support monitoring of vegetation response after fire \citep{perezcabello2021postfire}. Pe\~na-Molina et al.\ use a 33-year satellite time series to assess fire vulnerability, resilience, and recovery rates of Mediterranean pine forests, showing how long-term Earth-observation data can quantify recovery trajectories \citep{penamolina2024recovery}. Roche et al.\ analyze post-fire recovery dynamics and resilience of ecosystem-service capacity in Mediterranean-type ecosystems, connecting vegetation recovery with ecosystem-service restoration \citep{roche2024postfire}.

This literature is methodologically relevant because it demonstrates how ecological consequences can be measured through remote sensing, time-series analysis, and recovery indicators. For a loss-oriented wildfire thesis, such work suggests that future ecological-loss extensions could combine burn severity, vegetation recovery metrics, land-cover change, and ecosystem-service indicators to complement monetary loss targets.

\subsection{Biodiversity, Multifunctionality, and Long-Term Recovery}

A further ecological dimension concerns biodiversity and ecosystem multifunctionality. Laterza et al.\ show that wildfire severity and time since fire jointly shape biodiversity, ecosystem-service provision, and multifunctionality in Mediterranean forests \citep{laterza2026multifunctionality}. This is important because it indicates that ecological loss depends not only on whether a fire occurred, but also on severity, time since disturbance, and the recovery of multiple ecological functions. Busari et al.\ model wildfire effects on ecosystem services in two California watersheds and communities, illustrating how ecosystem-service impacts can be evaluated spatially after wildfire \citep{busari2025ecosystemservices}. Together, these studies strengthen the argument that ecological consequences require explicit indicators of severity, recovery, and ecosystem services rather than burned area alone.

For the present thesis, this expanded ecological literature suggests that ecological losses should be treated as a distinct but connected dimension of wildfire-loss assessment. Ecological damage may overlap with economic damage, but it is not reducible to it. A complete wildfire-loss framework should therefore incorporate ecological vulnerability, ecosystem services, recovery rates, and landscape-specific resilience alongside monetary impacts.

\section{LLMs, RAG, and Foundation Models in Wildfire Research}

\subsection{LLMs for Literature Synthesis and Decision Support}

Since 2024, wildfire research has begun to incorporate LLMs, multimodal systems, and foundation-model architectures. One use case is large-scale literature synthesis. Lin et al.\ use LLMs to analyze more than 60,000 wildfire research papers and identify geographic and thematic disparities in wildfire research \citep{lin2024llmresearch}. This is relevant because it shows that LLMs can help synthesize large research corpora, although this role is different from quantitative wildfire prediction.

Another use case is decision support. Xie et al.\ develop WildfireGPT, a domain-tailored LLM/RAG system designed to provide context-specific wildfire risk and resilience information \citep{xie2024wildfiregpt,xie2025rag}. Caron et al.\ propose a multi-target wildfire-risk prediction framework in which multiple predictive outputs are synthesized by LLMs into actionable reports for first responders \citep{caron2026multitarget}. Cheng et al.\ show that retrieval-informed prompting can improve agreement between LLM-generated rankings and expert-defined risk rankings in hierarchical wildfire risk evaluation \citep{cheng2026retrieval}. These studies suggest that LLMs are most useful when organizing, retrieving, and communicating complex wildfire knowledge rather than replacing domain-specific quantitative models.

\subsection{Vision-Language and Multimodal Systems}

Multimodal wildfire systems are also becoming more common. Ayanzadeh et al.\ introduce WildfireVLM, which combines satellite-based object detection with language-driven risk assessment and response recommendations \citep{ayanzadeh2026wildfirevlm}. Seidel et al.\ show how vision-language models can enhance UAV-based early wildfire detection by generating structured scene descriptions and improving situational awareness \citep{seidel2025advancing}. FireScope uses chain-of-thought-style conditioning to improve structured wildfire risk prediction and generalization \citep{markov2026firescope}. Xu et al.\ review generative AI for 2D and 3D wildfire spread prediction and argue that generative and multimodal models may become useful complements to physics-based and deep-learning workflows \citep{xu2025generativeai}.

\subsection{Limits of LLMs for Quantitative Wildfire Prediction}

Although these systems are promising, recent studies also make their limitations clear. Ramesh et al.\ compare WildfireGPT-style generative approaches with a domain-specific TabNet model for quantitative fire radiative power forecasting and show that the general-purpose generative system performs poorly for precise numerical prediction \citep{ramesh2025assessingwildfiregpt}. This suggests that LLMs and RAG systems should currently be viewed as augmentation layers for retrieval, explanation, and structured decision support, rather than standalone engines for quantitative wildfire forecasting or loss estimation.

This conclusion is particularly important for the present thesis. LLMs may help translate wildfire-risk and loss information into interpretable reports, but the underlying quantitative backbone still needs to come from domain-specific predictive and loss-estimation models.

\section{Summary of Related Work}

Table~\ref{tab:related_work_summary} summarizes the main categories of related work and their relevance to the present thesis.

\begin{table}[htbp]
\centering
\caption{Summary of related work categories and relevance to this thesis.}
\label{tab:related_work_summary}
\resizebox{\textwidth}{!}{%
\begin{tabular}{p{3.2cm}p{4.2cm}p{4.2cm}p{5.5cm}}
\hline
Category & Representative studies & Main focus & Relevance to thesis \\
\hline
Fire behavior and fire modeling & \citet{rothermel1972spread}, \citet{vanwagner1987fwi}, \citet{finney1998farsite}, \citet{scott2005fuelmodels} & Fuel, weather, topography, fire-spread systems & Provides physical basis for environmental predictors \\
Climate and fire regimes & \citet{westerling2006warming}, \citet{jolly2015climate}, \citet{abatzoglou2016anthropogenic}, \citet{williams2019observed} & Climate, aridity, fire-weather seasons, fire regimes & Justifies VPD, precipitation, temperature, and wind variables \\
Wildfire ML/DL prediction & \citet{jain2020review}, \citet{xu2024survey}, \citet{xu2025deeplearning}, \citet{ejaz2025mlpipeline}, \citet{caron2025ai} & Ignition, susceptibility, detection, spread, prediction pipelines & Provides modeling background and shows limitations of hazard-only prediction \\
Datasets and remote sensing & \citet{eidenshink2007mtbs}, \citet{giglio2016modis}, \citet{schroeder2014viirs}, \citet{xu2024sen2fire}, \citet{karlsson2025modisviirs} & Burn severity, active fire, satellite benchmarks & Supports wildfire data foundations and remote-sensing feature logic \\
Loss and risk assessment & \citet{asif2026supranational}, \citet{fema2025nri_technical}, \citet{huang2024suppression}, \citet{jia2025interpretable}, \citet{raveendran2025wildfireloss} & EAL, suppression costs, damage drivers, loss modeling & Directly supports the thesis target and loss-oriented framing \\
Cascading and indirect impacts & \citet{wang2021economicfootprint}, \citet{monge2023probabilistic}, \citet{ma2022supplychain}, \citet{qiu2025wildfiresmoke} & Economy-wide, supply-chain, smoke-health consequences & Broadens loss definition beyond direct burned area \\
Agricultural-hazard losses & \citet{lateef2025floodedcrop}, \citet{dang2024cropflood}, \citet{shi2021cropyield}, \citet{zhu2025ogallala} & Hazard-to-damage and recovery modeling & Provides methodological analogies for consequence-oriented modeling \\
Ecological losses and recovery & \citet{arrogantefunes2024ecological}, \citet{rocesdiaz2022ecosystemservices}, \citet{perezcabello2021postfire}, \citet{laterza2026multifunctionality} & Ecosystem vulnerability, ecosystem services, vegetation recovery, biodiversity & Strengthens the ecological-loss dimension of the thesis beyond monetary impacts \\
Spatial validation & \citet{roberts2017crossvalidation}, \citet{valavi2019blockcv}, \citet{meyer2021aoa}, \citet{wang2023spatialplus} & Blocked, grouped, and spatial validation & Justifies state-wise grouped validation in the experiments \\
LLMs and foundation models & \citet{xie2024wildfiregpt}, \citet{lin2024llmresearch}, \citet{cheng2026retrieval}, \citet{ramesh2025assessingwildfiregpt} & Retrieval, synthesis, decision support, limitations & Frames LLMs as support tools rather than replacements for quantitative models \\
\hline
\end{tabular}%
}
\end{table}

\section{Research Gap and Thesis Positioning}

The literature demonstrates substantial progress in wildfire occurrence modeling, remote-sensing-based prediction, spatiotemporal deep learning, ecological vulnerability mapping, ecosystem-service impact assessment, public risk indexing, and AI-assisted decision support. However, several gaps remain clear. First, much of the wildfire literature is still hazard-centric, with fewer studies translating predictive outputs into integrated economic and ecological loss estimates \citep{asif2026supranational,caron2025ai}. Second, although public risk datasets such as the FEMA NRI provide useful expected annual loss targets, more work is needed to connect these targets with environmental predictors and machine-learning validation frameworks \citep{fema2025nri_technical,fema2025nri_methodology}. Third, although adjacent natural-hazard fields demonstrate how remote sensing, damage indices, participatory mapping, and economic analysis can estimate consequences and recovery, equivalent integrated pipelines remain less common in wildfire studies \citep{lateef2025floodedcrop,dang2024cropflood,shi2021cropyield,zhu2025ogallala}. Fourth, ecological-loss studies increasingly examine ecosystem services, post-fire recovery, biodiversity, and multifunctionality, but these indicators are still rarely integrated with quantitative wildfire-risk and monetary-loss models \citep{rocesdiaz2022ecosystemservices,perezcabello2021postfire,roche2024postfire,laterza2026multifunctionality}. Fifth, spatial transferability remains an important open challenge because wildfire models can perform well under random validation but degrade when evaluated across new geographic regions \citep{roberts2017crossvalidation,valavi2019blockcv,meyer2021aoa,wang2023spatialplus}. Finally, while LLMs and foundation models are becoming useful for retrieval, synthesis, and interpretation, they are not yet reliable substitutes for quantitative wildfire-risk and wildfire-loss models \citep{xie2025rag,ramesh2025assessingwildfiregpt,cheng2026retrieval}.

Accordingly, the present thesis is positioned at the intersection of wildfire-risk prediction and wildfire-loss assessment. Its contribution lies not in treating wildfire hazard as an endpoint, but in viewing it as an input to a broader framework that considers what is exposed, how vulnerable it is, what environmental conditions shape wildfire likelihood and severity, and what economic and ecological consequences follow. The thesis responds to the need for integrated, geographically validated, consequence-oriented wildfire-loss modeling using accessible public datasets.
